% Generated draft: requires mathematical and source review.
\documentclass[11pt]{article}
\usepackage[margin=25mm]{geometry}
\usepackage{fontspec}
\setmainfont{Latin Modern Roman}
\usepackage{amsmath,amssymb,mathtools,unicode-math}
\setmathfont{Latin Modern Math}
\usepackage{enumitem,needspace,xurl,hyperref}
\setlength{\parindent}{0pt}
\setlength{\parskip}{5pt}
\setlength{\emergencystretch}{3em}
\newcommand{\N}{\mathbb{N}}
\newcommand{\Z}{\mathbb{Z}}
\newcommand{\Q}{\mathbb{Q}}
\newcommand{\R}{\mathbb{R}}
\newcommand{\C}{\mathbb{C}}
\newcommand{\F}{\mathbb{F}}
\newcommand{\A}{\mathbb{A}}
\newcommand{\PP}{\mathbb{P}}
\newcommand{\E}{\mathbb{E}}
\newcommand{\eps}{\varepsilon}
\newcommand{\dd}{\,\mathrm{d}}
\begin{document}
\section*{UnsolvedMath ID: 30001423}
\textbf{Source problem code:} catalogue entry
\subsection{Definitions and mathematical statement}
The supplied record concerns an unknown scalar Dirichlet-to-Neumann function \(\Lambda\) associated with a self-similar tree and displays the functional equation
\[
\Lambda(\omega)\cos\omega-\omega\sin\omega
=
\sum_{i=0}^{p-1}\frac{\mu_i}{\alpha_i}
\left(\cos\omega+\frac{\sin\omega}{\omega}\Lambda(\alpha_i\omega)\right).
\]
Here \(\omega\) is the frequency variable, and \(p,\mu_i,\alpha_i\) are parameters whose admissible values and geometric meanings are not specified in the supplied statement. The frequency domain, scalar field, tree construction, boundary conditions, and regularity requirements on \(\Lambda\) are also unspecified. The displayed expressions require nonzero denominators wherever they are evaluated; no interpretation at \(\omega=0\) is supplied.

The requested assertion is that there are at most two ``homeomorphic solutions'' and that \(\Lambda(0)\) uniquely determines a solution. To express only its quantifier structure, let \(\mathcal S\) stand for the intended, presently unspecified solution class for fixed admissible tree data and parameters. Provided evaluation at zero is defined for every member, the proposed conclusions would be
\[
|\mathcal S|\leq 2,
\qquad
\forall\Lambda_1,\Lambda_2\in\mathcal S,\quad
\Lambda_1(0)=\Lambda_2(0)\Longrightarrow\Lambda_1=\Lambda_2.
\]
The task is to prove or disprove these conclusions once the missing specifications are recovered. This notation does not define the phrase ``homeomorphic solutions'': the source gives neither spaces between which the solutions should be homeomorphisms nor another precise meaning for that phrase. Consequently, these formulas record the intended claim but do not constitute a recovered, well-posed conjecture. The supplied statement-verification record marks the formulation as unrecoverable and supplies no clean replacement.
\subsection{Short English statement}
The source asks whether the displayed self-similar-tree Dirichlet-to-Neumann equation has at most two solutions in a class called ``homeomorphic,'' with each solution uniquely determined by its value at zero. That class, the parameter assumptions, the frequency domain, and the interpretation at zero are missing, so a precise conjecture cannot yet be recovered.
\subsection{Sources}
\begin{itemize}
\item Local UnsolvedMath registry, supplied complete record with permanent ID \(30001423\), source code \texttt{catalogue entry}, titled ``Incomplete Self-Similar-Tree DtN Uniqueness Claim.'' The record belongs to the Oberwolfach Reports Open Problems collection and assigns the proposed year 2010; no individual proposer is supplied.
\item \emph{Computational Electromagnetism and Acoustics} (2010), Oberwolfach Reports, DOI \texttt{10.4171/owr/2010/10}, pp.\ 415--519. Supplied source report: \url{https://doi.org/10.4171/owr/2010/10}. The background describes Oberwolfach Reports as documenting mathematical workshops and their open-problem discussions.
\item Joly--Kachanovska--Semin, 2019, supplied literature reference: \url{https://doi.org/10.3934/nhm.2019010}.
\item Oberwolfach Report 10/2010, supplied literature reference: \url{https://doi.org/10.4171/owr/2010/10}.
\item Literature assessment checked at \texttt{2026-08-21}: the supplied curation labels the problem ``partially solved'' and states that Joly, Kachanovska, and Semin prove uniqueness of the relevant DtN functional equation under appropriate assumptions and develop the fractal-tree theory, while the full homeomorphism and at-most-two formulation is not established in the cited work. This is supplied dated curation derived from the OWR assessment, not an independently verified literature assessment; its status and cited claims require independent verification.
\item The supplied statement-verification field marks the statement as ``unrecoverable'' because ``homeomorphic solutions'' is undefined for the displayed scalar function as written. No separate clean statement is promoted; the original record is retained publicly for provenance. The supplied record was updated at \texttt{2026-09-16T12:06:53Z}.
\end{itemize}
\end{document}
